% ---------------------------------------------------------------
%  Section I -- Introduction
% ---------------------------------------------------------------

Urban mobility in rapidly expanding metropolitan areas is shaped by three
interacting structural problems: peak-hour congestion that imposes measurable
economic and environmental costs~\cite{schrank2019tti}, fragmented data
held simultaneously by multiple operators and public agencies, and
regulatory constraints that make centralizing raw personal traces infeasible
in many jurisdictions.
Effective real-time prediction and optimization require integrating
road-network topology, sensor telemetry, transit schedules, atmospheric
conditions, and large-scale event calendars, yet no single operator
typically controls all of these inputs.

Baku, the capital of Azerbaijan and the country's largest urban region,
exemplifies this challenge acutely.
The city has substantially expanded its road network over the past decade
and is investing in an extended metro system and bus-rapid-transit corridors,
yet public, machine-readable GTFS feeds were not available to the authors
during preparation of the baseline release.
Planners and travelers simultaneously need short-horizon congestion forecasts
that can inform signal control and navigation, multimodal journey plans
accounting for current service disruptions, equity-aware accessibility
metrics that identify underserved districts, and incident detection that
distinguishes genuine demand shocks from sensor faults.
In the absence of open data feeds, any deployed platform must communicate
its data-source status explicitly rather than silently degrading or
substituting synthetic data.

\textit{IRIDIUM} is an open-source platform that addresses this operational
and architectural gap through four pillars.

\noindent\textbf{Dynamic digital twin.}
A directed attributed mobility graph $G=(V,E,W)$ is assembled at each
request from all configured real sources.
Nodes represent junctions, road segments, and transit stops; edges carry
mode-specific travel time, monetary cost, carbon proxy, and dynamic state
attributes.
When a source is absent or unconfigured, the assembler returns an empty but
validly typed payload together with an explicit status code, never a
synthetic substitute.

\noindent\textbf{Modular analytics suite.}
A congestion forecaster targets the multi-step prediction
$\hat{\mathbf{Y}}_{t+1:t+H} = f_\theta(\mathbf{X}_{t-T+1:t}, G)$
via a spatio-temporal graph neural network (ST-GNN) extension point with a
heuristic fallback for the pre-model baseline.
A constrained shortest-path router minimizes the weighted functional
$J(r)=\alpha T(r)+\beta C(r)+\gamma E(r)+\delta P(r)$.
A district-level Mobility Equity Score $\MES_d=\sum_m w_m z_{d,m}$
supports spatial justice analysis.
A composite anomaly severity score
$S_t = \lambda_1|z_t|+\lambda_2 e_t+\lambda_3 w_t$ fuses speed deviations
with event and weather signals.

\noindent\textbf{Federated learning pathway.}
Operator nodes train local models on their private datasets and contribute
parameter updates to a central aggregator following the FedAvg rule
$\theta^{(r+1)} = \sum_k p_k\theta_k^{(r)}$~\cite{mcmahan2017fedavg}.
The pathway supports optional $(\varepsilon,\delta)$-differential privacy
noise calibration to protect against gradient inversion
attacks~\cite{kairouz2021advances}.
No federated training experiments are reported in this paper; the FL
component is fully architected and awaits deployment of local training nodes.

\noindent\textbf{Provenance-and-status protocol.}
Every API response carries a machine-readable \texttt{data\_status} tag
and a source provenance list $\Pi = [\pi_{s_1},\ldots,\pi_{s_K}]$
(Eq.~\ref{eq:provenance}).
This design ensures that downstream consumers can always determine whether
a result is grounded in real observations and enables automated data-quality
monitoring as new sources are activated.

\subsection*{Gap Addressed}

Prior work on urban mobility prediction largely relies on either simulated
traffic generation or on proprietary datasets that are unavailable to
independent replication~\cite{pan2019urban,yao2019revisiting,wang2020traffic}.
Platforms built for cities in North America or Western Europe
(METR-LA, PeMS, London TfL) benefit from mature open-data
ecosystems~\cite{li2018dcrnn,wu2019graphwavenet} that do not yet exist
in the South Caucasus region.
IRIDIUM establishes a \textit{real-data stack} in which every source is
documented with acquisition method, license, refresh cadence, and fallback
behavior, enabling honest evaluation under partial data availability and
incremental deployment as operator agreements and feed access become
available.

\subsection*{Contributions}

The primary contributions of this paper are as follows.
\begin{enumerate}
  \item A precise, self-consistent mathematical formalization of the
        mobility graph, the multi-objective routing functional,
        the ST-GNN forecasting target, the Mobility Equity Score,
        the composite anomaly severity score, and the FedAvg aggregation
        step, each consistent with the implemented codebase
        (Section~\ref{sec:formulation}).
  \item A layered system architecture specification that makes the
        provenance-and-status protocol a first-class citizen of
        every API response, supported by a formal description of the
        five-layer platform design and the state-assembly algorithm
        (Tables~\ref{tab:layers}--\ref{tab:pipeline},
        Section~\ref{sec:formulation}).
  \item A real-data ingestion stack covering OpenStreetMap/PostGIS,
        Open-Meteo, and a provider-agnostic traffic and transit
        abstraction layer with quantified data-access status for all
        eight registered sources
        (Section~\ref{sec:data}).
  \item A baseline evaluation covering API semantic correctness,
        end-to-end response latency with fitted log-normal parameters,
        and a real 72-hour weather time-series integration for Baku
        and Quba, including a quantification of the persistent
        inter-city temperature differential
        (Sections~\ref{sec:setup}--\ref{sec:results}).
  \item A candid, mathematically grounded accounting of all unresolved
        operator dependencies, data-access gaps, and evaluation
        limitations (Section~\ref{sec:limitations}).
\end{enumerate}

The IRIDIUM codebase, configuration files, and dataset documentation
are publicly available at
\url{https://github.com/iridium-oss/iridium}.

The remainder of this paper is organized as follows.
Section~\ref{sec:related} reviews related work.
Section~\ref{sec:formulation} formalizes the system model.
Section~\ref{sec:data} describes the data sources.
Section~\ref{sec:methods} details the implemented pipeline.
Section~\ref{sec:setup} explains the experimental setup.
Section~\ref{sec:results} presents results.
Section~\ref{sec:discussion} discusses implications.
Section~\ref{sec:limitations} addresses limitations, ethics, and privacy.
Section~\ref{sec:conclusion} concludes.
